\begin{frame}[allowframebreaks, t]
    \frametitle{De generadores a ecuaciones implícitas}

    \begin{proposition}
        Sean $v_1,\ldots,v_r\in k^n$ y $A=(v_1\mid\cdots\mid v_r)$.
        Las ecuaciones implícitas de $W=\langle v_1,\ldots,v_r\rangle$
        se obtienen escalonando la matriz $(A\mid u)$, donde
        $u=(u_1,\ldots,u_n)$ es simbólico: cada fila en que la parte
        de coeficientes se anule da una condición lineal
        $a_1u_1+\cdots+a_nu_n=0$, y $W$ es exactamente el conjunto de
        vectores que satisfacen todas esas condiciones.
    \end{proposition}

    \begin{proof}
        Un vector $u\in k^n$ pertenece a $W$ si y solo si el sistema
        $Ax=u$ es compatible. Al escalonar $(A\mid u)$ con $u$ simbólico,
        las filas en que la parte de coeficientes se anula expresan
        exactamente las condiciones de compatibilidad sobre $u_1,\ldots,u_n$.
    \end{proof}

    En particular, si $\dim W=s$, hay $n-s$ ecuaciones implícitas
    independientes.
\end{frame}


\begin{frame}[allowframebreaks, t]
    \frametitle{Ejemplo: ecuaciones implícitas del plano}

    \begin{example}
        Hallemos las ecuaciones implícitas de
        $W=\langle(2,1,0),(-1,0,1)\rangle\subseteq\R^3$.
    \end{example}

    Escalonamos la matriz ampliada $(A\mid u)$:
    \[
        \begin{pmatrix}
            2&-1&u_1\\1&0&u_2\\0&1&u_3
        \end{pmatrix}
        \xrightarrow{\;F_1\leftrightarrow F_2\;}
        \begin{pmatrix}
            1&0&u_2\\2&-1&u_1\\0&1&u_3
        \end{pmatrix}
        \xrightarrow{\;F_2-2F_1\;}
        \begin{pmatrix}
            1&0&u_2\\0&-1&u_1-2u_2\\0&1&u_3
        \end{pmatrix}
    \]
    \[
        \xrightarrow{\;F_3+F_2\;}
        \begin{pmatrix}
            1&0&u_2\\0&-1&u_1-2u_2\\0&0&u_1-2u_2+u_3
        \end{pmatrix}.
    \]
    La tercera fila da la condición $u_1-2u_2+u_3=0$. Por tanto,
    \[
        W\colon u_1-2u_2+u_3=0,
    \]
    que coincide con la ecuación del plano del ejemplo anterior.
\end{frame}
